\section{The Fundamental Counting Principle}
The probability of winning top prize in the lottery is about the same as the probability of being struck by lightening. There are million possible number combinations in the lottery games, and only one way of winning the grand prize.

%\subsection{The fundamental counting principle with two groups of items}
	We can generalize the idea to any two groups of items with the {\bf Fundamental Counting Principle.}

\begin{tcolorbox}[enhanced,width=6.65in,center upper,size=fbox,drop shadow southwest,sharp corners]

\begin{tem}[The Fundamental Counting Principle] with two groups of items: If you can select one item from a group of {\bf M } items and another item from a group of {\bf N} items, then the total number of possible two-item combinations is {\bf M $\cdot$ N.}
\end{tem}
\end{tcolorbox}	

\begin{exe} The Greasy Spoon Restaurant offers 6 appetizers and 14 main courses. In how many ways can a person order a two-course meal?
\end{exe}
%\begin{sol}
%Choose from one of 6 appetizers and one of 14 main courses, the number of two-course meals is
%\[6\cdot 14=84.\]
%A person can order a two-course meal in 84 different ways.
%\end{sol}
\vfill

\begin{ex} A restaurant offers 10 appetizers and 15 main courses. In how many ways can a person order a two-course meal?
\end{ex}
\vfill

\begin{ex} This is a semester that you will take your required psychology and social sciences courses. Because you decide to register early, there are 15 section of psychology from which you can choose. Moreover, there are 9 sections of social science that are available at times that do not conflict with those of psychology. In how many ways you create two-course schedules that satisfy the psychology-social science requirement?
\end{ex}
\vfill
\pagebreak

%\begin{sol}
%The number of ways that you can satisfy the requirement is found by multiplying the number of choices for each course. You can choose your psychology course from 15 sections and your social science course from 9 sections. For both courses you have
%\[15 \cdot 9, \quad or \quad 135\]
%Choices. Thus, you can satisfy the psychology-social science requirement in 135 ways.
%\end{sol}
%\vfill

\begin{tcolorbox}[enhanced,width=6.5in,center upper,size=fbox,drop shadow southwest,sharp corners]

\begin{tem}[The Fundamental Counting Principle]  
%{\bf with more than two groups of items}: 
The number of ways in which a series of successive events can occur is determined {\bf by multiplying} the number of ways each individual event can happen.

%The number of ways in which a series of successive things can occur is found {\bf by multiplying} the number of ways in which each thing can occur.
\end{tem}
\end{tcolorbox}


\begin{exe} 
A pizza can be ordered with two choices of size (medium or large), three choices of crust (thin, thick or regular), and five choices of toppings (ground beef, sausage, pepperoni, bacon and mushrooms). How many different one-topping pizzas can be ordered?
\end{exe}
%\begin{sol}
%We use the Fundamental Counting Principle to find the number of different one-topping pizzas.
%\[2 \cdot 3 \cdot 5 = 30.\]
%Thus, there are 30 different one-topping pizzas that we can order.
%\end{sol}
\vfill
%\pagebreak

\begin{exe}\label{4}
Car manufactures are now experimenting with lightweight three-wheel cars, designed for one person, and considered ideal for city driving. Intrigued? Suppose you could order such a car with a choice of 9 possible colors, with or without air conditioning, electric or gas powered, and with or without an on board computer. In how many ways can this car be ordered with regard to these options?
\end{exe}
%\begin{sol}
%This situation involves making choices with four groups of items.
%\[Color(9\ choices) \quad Air-conditioning(2)  \quad Power(2)  \quad Computer(2)  \]
%We use the Fundamental Counting Principle to find the number of ordering options.
%\[9 \cdot 2\cdot 2 \cdot 2=72.\]
%Thus, the car can be ordered in 72 different ways.
%\end{sol}
\vfill

\begin{ex} 
The car in Example \ref{4} is now available in 10 possible colors. The option involving air conditioning, power, and on board computer still apply. Furthermore, the car is available with or without a global positioning system. In how many ways can this car be ordered in terms of these options?
\end{ex}
\vfill
\pagebreak

\begin{exe}
You are taking a multiple-choice test that has ten questions. Each of the questions has four answer choices, with one correct answer per question. If you select one of these four choices for each question and leave nothing blank, in how many ways can you answer the questions?
\end{exe}

%\begin{solution}
{\bf Solution}: 
This situation involves making choices with ten questions.
\[ Question\ 1 (4 \ choices) \quad Question\ 2 \quad Question\ 3 \quad \cdots \quad  Question \ 10.\]
We use the Fundamental Counting Principle to find the number of ways that you can answer the questions on the test. Multiply the number of choices, 4, for each of the ten questions.
\[4 \cdot 4 \cdot 4 \cdot 4 \cdot 4 \cdot 4 \cdot 4 \cdot 4 \cdot 4 \cdot 4 =4^{10}=1,048,576.\]
Thus, you can answer the questions in 1,048,576 different ways.
%\end{solution}

\begin{ex}
You are taking a multiple-choice test that has six questions. Each of the questions has three answer choices, with one correct answer per question. If you select one of these three choices for each question and leave nothing blank, in how many ways can you answer the questions?
\end{ex}
\vfill


\begin{ex}  
Telephone numbers in the United States begin with three-digit area codes followed by seven-digit local telephone numbers. Area codes and local telephone numbers cannot begin with 0 or 1. How many different telephone numbers are possible?
\end{ex}
%{\bf Solution.} 
%Here are the choices for each of the ten groups of items:
%\[Area\ Code \quad \quad Local\ telephone \ Number\]
%\[508 \quad \quad \quad \quad 929 \quad \quad 8000 \quad \quad  \]
%We use the Fundamental Counting Principle to determine the number of different telephone numbers that are possible. The total number of telephone numbers possible is 
%\[ 8 \cdot 10 \cdot 10 \cdot 8 \cdot 10 \cdot 10\cdot 10\cdot 10\cdot 10\cdot 10=6,400,000,000.\]
%There are six billion four hundred million different telephone numbers that are possible.
\vfill
\pagebreak

\begin{ex} 
A popular type of pen comes in red, blue, or black ink.  The writing tip varies from extra bold, bold, regular, fine, or micro.  How many different choices of pens do you have with this type of pen?
\end{ex}
\vfill

\begin{ex}
A wedding caterer gives you three choices for the main course, six starter choices and five options for dessert. How many different meals (made up of starter, dinner and dessert) are there?
\end{ex}
\vfill

\begin{ex}
You take a survey with five ``yes'' or ``no'' answers. How many different ways could you complete the survey?
\end{ex}
\vfill
\begin{ex}
A company puts a code on each different product they sell. The code is made up of 3 numbers and 2 letters. How many different codes are possible?
\end{ex}
\vfill
\pagebreak

%\section{Factorials}
%	The product 
%	\[5 \cdot 4\cdot 3 \cdot 2 \cdot 1 \]
%	is given a special name and symbol. It is called 5 {\bf factorial}, and written as {\bf 5!} Thus,
%	\[5!=5 \cdot 4\cdot 3 \cdot 2 \cdot 1.\]
%	In general, 
If $n$ is a positive integer, then $n!$ ({\bf $n$ factorial}) is the product of all positive integers from $n$ down through 1. For example,
	\[
	\begin{array}{ccl}
	1! & = & 1\\
	2! & = & 2 \cdot 1 = 2 \\
	3! & = & 3 \cdot 2 \cdot 1  = 6\\
	4! & = & 4\cdot 3 \cdot 2 \cdot 1 = 24\\
	5! & = & 5 \cdot 4\cdot 3 \cdot 2 \cdot 1=120\\
%	6! & = & 6 \cdot 5 \cdot 4\cdot 3 \cdot 2 \cdot 1=720\\
%	7! & = & 7 \cdot 6 \cdot 5 \cdot 4\cdot 3 \cdot 2 %\cdot 1=5040.
	\end{array}
	\]	
	
\begin{tcolorbox}[enhanced,width=6.5in,center upper,size=fbox,drop shadow southwest,sharp corners]
{\bf Factorials}: %If $n$ is a positive integer, the notation $n!$ ($n$ factorial) is the product of all positive integers from $n$ down through 1.
$$\displaystyle n!=n(n-1)(n-2)(n-3)\cdots (3)(2)(1)$$   
0! (zero factorial), by definition, is 1 $0!=1.$	
\end{tcolorbox}

\begin{exe}
Evaluate the following factorial expressions without using the factorial key on your calculator.
\begin{multicols}{1}
\begin{enumerate}
	\item $\dfrac{8!}{5!}$\vfill
	\item $\dfrac{21!}{24!}$\vfill
	\item $\dfrac{500!}{499!}$\vfill
	\item $\dfrac{11!}{7!}$\vfill
	\item $\dfrac{13!}{16!}$\vfill
	\item $\dfrac{99!}{100!}$\vfill
	\item $\dfrac{7!}{(5-1)!}$\vfill
	\item $\dfrac{7!}{3!(4-4)!}$\vfill
\end{enumerate}
\end{multicols}
\end{exe}\vfill
\pagebreak

\begin{tcolorbox}[enhanced,width=6.5in,center upper,size=fbox,drop shadow southwest,sharp corners]

{\bf Permutations of duplicate items}: 
The number of permutations of $n$ items, where $p$ items are identical, $q$ items are identical, $r$ items are identical, $s$ items are identical, and so on, is given by
$$\displaystyle\frac{n!}{(p!\ q!\ r!\ s! \cdots)}.$$  
\end{tcolorbox}

\begin{exe}
In how many distinct ways can the letters of the word\\ ``MISSISSIPPI'' be arranged?
\end{exe}\vfill
%\pagebreak 
%\begin{sol}
%{\bf Solution}: The word contains 11 letters where four Is are identical, four Ss are identical and 2 Ps are identical. The number of distinct permutations is
%\[\displaystyle \frac{n!}{p!q!r!}=\frac{11!}{(4!4!2!)}=34,650.\]
%\end{sol}

\begin{ex}
In how many distinct ways can the letters of the word ``BANANA'' be arranged?
\end{ex}
\vfill
\begin{ex}
In how many distinct ways can the letters of the word\\ ``MASSACHUSETTS'' be arranged?
\end{ex}
\vfill
%\begin{ex}
%In how many distinct ways can the letters of the word %``TENNESSEE'' be arranged?
%\end{ex}
%\vfill
\begin{ex}
There are seventeen balloons: 3 blue, 5 red, 2 green, 3 yellow, and 4 orange. In how many distinct ways can the balloons be arranged?
\end{ex}
\vfill
\pagebreak
